% --- Archon physics formalization source begin ---
% archon:physics
% archon:covers PhyXMiniProblems/problem_phyx_mini_0041.lean
% archon:source-report reports/phyx_mini/problem_phyx_mini_0041.source.json

\chapter{Physics problem phyx\_mini\_0041}
\label{ch:PhyXMiniProblems_problem_phyx_mini_0041}

\paragraph{Problem source.}
\#\# Physical scenario

A cylindrical glass rod (figure) has index of refraction 1.52. It is surrounded by water that has index of refraction of 1.33. One end is ground to a hemispherical surface with radius \textbackslash{}( R = 2.00 \textbackslash{} \textbackslash{}mathrm\{cm\} \textbackslash{}). A small object is placed on the axis of the rod, 8.00 \textbackslash{} \textbackslash{}mathrm\{cm\} to the left of the vertex.

\#\# Question

Find the lateral magnification.

\#\# Answer choices

- A: A: +0.856

- B: B: +2.33

- C: C: +0.995

- D: D: +1.814

\#\# Figure caption (auxiliary; use the image as primary evidence)

The image depicts a ray diagram illustrating light refraction between two media. 

1. **Background and Media**:
   - The left side is labeled as having a refractive index \textbackslash{}( n\_a = 1.33 \textbackslash{}), indicating it represents water.
   - The right side has a refractive index \textbackslash{}( n\_b = 1.52 \textbackslash{}), suggesting a different medium, possibly glass or another material.

2. **Objects and Points**:
   - Point \textbackslash{}( P' \textbackslash{}) is marked on the left side (pink).
   - Point \textbackslash{}( P \textbackslash{}) is at the boundary between the media (blue).
   - Point \textbackslash{}( C \textbackslash{}) is in the right medium (black).

3. **Lines and Arrows**:
   - Arrows originating from point \textbackslash{}( P \textbackslash{}) to the right are refracted light rays passing through to the region \textbackslash{}( n\_b \textbackslash{}).
   - Dashed lines from \textbackslash{}( P' \textbackslash{}) are aimed towards \textbackslash{}( P \textbackslash{}), representing incident light paths.
   - Solid horizontal lines represent distances, with the distance \textbackslash{}( s = 8.00 \textbackslash{}text\{ cm\} \textbackslash{}) marked between point \textbackslash{}( P \textbackslash{}) and the right side.

4. **Measurements**:
   - \textbackslash{}( s' \textbackslash{}) denotes the distance from \textbackslash{}( P' \textbackslash{}) onwards, with the arrow indicating its measurement direction.

Overall, the diagram explains the refraction of light as it passes from water (left) to a denser medium (right), illustrating concepts such as incident rays and refracted rays based on Snell's law.

\#\# Dataset metadata

Geometrical Optics; Physical Model Grounding Reasoning, Multi-Formula Reasoning

\paragraph{Recorded answer/context.}
B: B: +2.33

\paragraph{Figure/image path.}
/root/proposal\_for\_physic/science-mango/phyx\_mini\_run/phyx\_data/test\_image/41.png

\paragraph{Formalization target.}
create a compiling Lean file with sorry bodies at `PhyXMiniProblems/problem\_phyx\_mini\_0041.lean`. The Lean declarations must preserve the physical quantities, dimensions or dimensional roles, figure labels, governing-law hypotheses, and final relation expressed by this problem.
Use Mathlib/Physlib names found through LeanExplore where available. If a physics API is missing, introduce faithful local abstractions rather than scalar placeholder aliases.

\begin{theorem}[Physics formalization target]
\label{thm:physics:phyx_mini_0041:target}
\lean{PhyXMiniProblems.ProblemPhyXMini0041.problem_phyx_mini_0041}
\uses{def:physics:phyx-mini-0041:phyxminiproblems-problemphyxmini0041-cylindricalglassrodsetup, def:physics:phyx-mini-0041:phyxminiproblems-problemphyxmini0041-hasstatedrefractiveindices, def:physics:phyx-mini-0041:phyxminiproblems-problemphyxmini0041-hasphysicalsigns, def:physics:phyx-mini-0041:phyxminiproblems-problemphyxmini0041-matchesfigurereadout, def:physics:phyx-mini-0041:phyxminiproblems-problemphyxmini0041-obeysexactsnelllawnearopticalaxis, def:physics:phyx-mini-0041:phyxminiproblems-problemphyxmini0041-hasfirstordersphericalraygeometry, def:physics:phyx-mini-0041:phyxminiproblems-problemphyxmini0041-mapsopticalaxistoitself, def:physics:phyx-mini-0041:phyxminiproblems-problemphyxmini0041-obeysexactchiefraysnelllawnearaxis, def:physics:phyx-mini-0041:phyxminiproblems-problemphyxmini0041-haslateralmagnificationatopticalaxis, def:physics:phyx-mini-0041:phyxminiproblems-problemphyxmini0041-matchesanswerchoice}
The assigned autoformalize agent should translate this physics problem into Lean declarations in the covered file, with theorem and lemma proof bodies written as `by sorry`.
\end{theorem}
\begin{proof}
This is an autoformalization task, not a proof task. Produce statements that can later be proved by the physics prover without weakening the physical model.
\end{proof}

% --- Archon declaration topology begin ---
\section{Declaration topology}
\begin{definition}[Dim Length]
  \label{def:physics:phyx-mini-0041:phyxminiproblems-problemphyxmini0041-dimlength}
  \lean{PhyXMiniProblems.ProblemPhyXMini0041.DimLength}
  A signed, unit-independent physical length
\end{definition}
\begin{proof}
  This is the declared model definition of Dim Length.
\end{proof}
\begin{definition}[centimeter Unit Choices]
  \label{def:physics:phyx-mini-0041:phyxminiproblems-problemphyxmini0041-centimeterunitchoices}
  \lean{PhyXMiniProblems.ProblemPhyXMini0041.centimeterUnitChoices}
  Unit choices in which lengths are read in centimeters and other units are SI
\end{definition}
\begin{proof}
  This is the declared model definition of centimeter Unit Choices.
\end{proof}
\begin{definition}[length In Centimeters]
  \label{def:physics:phyx-mini-0041:phyxminiproblems-problemphyxmini0041-lengthincentimeters}
  \lean{PhyXMiniProblems.ProblemPhyXMini0041.lengthInCentimeters}
  \uses{def:physics:phyx-mini-0041:phyxminiproblems-problemphyxmini0041-dimlength, def:physics:phyx-mini-0041:phyxminiproblems-problemphyxmini0041-centimeterunitchoices}
  The signed scalar readout of a physical length in centimeters
\end{definition}
\begin{proof}
  This is the declared model definition of length In Centimeters.
\end{proof}
\begin{definition}[Optical Medium]
  \label{def:physics:phyx-mini-0041:phyxminiproblems-problemphyxmini0041-opticalmedium}
  \lean{PhyXMiniProblems.ProblemPhyXMini0041.OpticalMedium}
  The optical media on the incident and transmitted sides of the surface
\end{definition}
\begin{proof}
  This is the declared model definition of Optical Medium.
\end{proof}
\begin{definition}[Figure Point]
  \label{def:physics:phyx-mini-0041:phyxminiproblems-problemphyxmini0041-figurepoint}
  \lean{PhyXMiniProblems.ProblemPhyXMini0041.FigurePoint}
  Axial point labels appearing in the primary ray diagram
\end{definition}
\begin{proof}
  This is the declared model definition of Figure Point.
\end{proof}
\begin{definition}[Cylindrical Glass Rod Setup]
  \label{def:physics:phyx-mini-0041:phyxminiproblems-problemphyxmini0041-cylindricalglassrodsetup}
  \lean{PhyXMiniProblems.ProblemPhyXMini0041.CylindricalGlassRodSetup}
  \uses{def:physics:phyx-mini-0041:phyxminiproblems-problemphyxmini0041-dimlength, def:physics:phyx-mini-0041:phyxminiproblems-problemphyxmini0041-opticalmedium, def:physics:phyx-mini-0041:phyxminiproblems-problemphyxmini0041-figurepoint}
  The physical setup and its local ray observables. The argument of each angle profile is the signed height, in centimeters, at which a ray meets the surface. Angle values are dimensionless radian readouts. The image-height map takes a small signed object height in centimeters to the corresponding signed paraxial image height in centimeters. Its derivative at zero is therefore dimensionless. No field fixes the requested numerical magnification or the virtual-image distance.
\end{definition}
\begin{proof}
  This is the declared model definition of Cylindrical Glass Rod Setup.
\end{proof}
\begin{definition}[axial Coordinate Centimeters]
  \label{def:physics:phyx-mini-0041:phyxminiproblems-problemphyxmini0041-axialcoordinatecentimeters}
  \lean{PhyXMiniProblems.ProblemPhyXMini0041.axialCoordinateCentimeters}
  \uses{def:physics:phyx-mini-0041:phyxminiproblems-problemphyxmini0041-lengthincentimeters, def:physics:phyx-mini-0041:phyxminiproblems-problemphyxmini0041-figurepoint, def:physics:phyx-mini-0041:phyxminiproblems-problemphyxmini0041-cylindricalglassrodsetup}
  The centimeter coordinate of a labelled point on the optical axis
\end{definition}
\begin{proof}
  This is the declared model definition of axial Coordinate Centimeters.
\end{proof}
\begin{definition}[object Distance Centimeters]
  \label{def:physics:phyx-mini-0041:phyxminiproblems-problemphyxmini0041-objectdistancecentimeters}
  \lean{PhyXMiniProblems.ProblemPhyXMini0041.objectDistanceCentimeters}
  \uses{def:physics:phyx-mini-0041:phyxminiproblems-problemphyxmini0041-cylindricalglassrodsetup, def:physics:phyx-mini-0041:phyxminiproblems-problemphyxmini0041-axialcoordinatecentimeters}
  Positive object distance s, measured leftward from the vertex to P
\end{definition}
\begin{proof}
  This is the declared model definition of object Distance Centimeters.
\end{proof}
\begin{definition}[signed Image Distance Centimeters]
  \label{def:physics:phyx-mini-0041:phyxminiproblems-problemphyxmini0041-signedimagedistancecentimeters}
  \lean{PhyXMiniProblems.ProblemPhyXMini0041.signedImageDistanceCentimeters}
  \uses{def:physics:phyx-mini-0041:phyxminiproblems-problemphyxmini0041-cylindricalglassrodsetup, def:physics:phyx-mini-0041:phyxminiproblems-problemphyxmini0041-axialcoordinatecentimeters}
  Signed paraxial image distance s', negative for P' on the water side
\end{definition}
\begin{proof}
  This is the declared model definition of signed Image Distance Centimeters.
\end{proof}
\begin{definition}[signed Radius Centimeters]
  \label{def:physics:phyx-mini-0041:phyxminiproblems-problemphyxmini0041-signedradiuscentimeters}
  \lean{PhyXMiniProblems.ProblemPhyXMini0041.signedRadiusCentimeters}
  \uses{def:physics:phyx-mini-0041:phyxminiproblems-problemphyxmini0041-cylindricalglassrodsetup, def:physics:phyx-mini-0041:phyxminiproblems-problemphyxmini0041-axialcoordinatecentimeters}
  Signed vertex-to-center distance, positive toward the glass rod
\end{definition}
\begin{proof}
  This is the declared model definition of signed Radius Centimeters.
\end{proof}
\begin{definition}[Has Stated Refractive Indices]
  \label{def:physics:phyx-mini-0041:phyxminiproblems-problemphyxmini0041-hasstatedrefractiveindices}
  \lean{PhyXMiniProblems.ProblemPhyXMini0041.HasStatedRefractiveIndices}
  \uses{def:physics:phyx-mini-0041:phyxminiproblems-problemphyxmini0041-cylindricalglassrodsetup}
  The dimensionless refractive-index readouts stated in the problem
\end{definition}
\begin{proof}
  This is the declared model definition of Has Stated Refractive Indices.
\end{proof}
\begin{definition}[Has Physical Signs]
  \label{def:physics:phyx-mini-0041:phyxminiproblems-problemphyxmini0041-hasphysicalsigns}
  \lean{PhyXMiniProblems.ProblemPhyXMini0041.HasPhysicalSigns}
  \uses{def:physics:phyx-mini-0041:phyxminiproblems-problemphyxmini0041-lengthincentimeters, def:physics:phyx-mini-0041:phyxminiproblems-problemphyxmini0041-cylindricalglassrodsetup, def:physics:phyx-mini-0041:phyxminiproblems-problemphyxmini0041-objectdistancecentimeters}
  Positivity and nondegeneracy of the physical data used as denominators
\end{definition}
\begin{proof}
  This is the declared model definition of Has Physical Signs.
\end{proof}
\begin{definition}[Matches Figure Readout]
  \label{def:physics:phyx-mini-0041:phyxminiproblems-problemphyxmini0041-matchesfigurereadout}
  \lean{PhyXMiniProblems.ProblemPhyXMini0041.MatchesFigureReadout}
  \uses{def:physics:phyx-mini-0041:phyxminiproblems-problemphyxmini0041-lengthincentimeters, def:physics:phyx-mini-0041:phyxminiproblems-problemphyxmini0041-cylindricalglassrodsetup, def:physics:phyx-mini-0041:phyxminiproblems-problemphyxmini0041-axialcoordinatecentimeters, def:physics:phyx-mini-0041:phyxminiproblems-problemphyxmini0041-objectdistancecentimeters, def:physics:phyx-mini-0041:phyxminiproblems-problemphyxmini0041-signedradiuscentimeters}
  Problem and figure readouts. The object P is 8.00 cm left of the vertex, the center C is one radius to its right, and that radius is 2.00 cm. The dashed backward extensions place the paraxial virtual image P' farther left than P. No numerical image distance or magnification is read from the figure
\end{definition}
\begin{proof}
  This is the declared model definition of Matches Figure Readout.
\end{proof}
\begin{definition}[Obeys Exact Snell Law Near Optical Axis]
  \label{def:physics:phyx-mini-0041:phyxminiproblems-problemphyxmini0041-obeysexactsnelllawnearopticalaxis}
  \lean{PhyXMiniProblems.ProblemPhyXMini0041.ObeysExactSnellLawNearOpticalAxis}
  \uses{def:physics:phyx-mini-0041:phyxminiproblems-problemphyxmini0041-cylindricalglassrodsetup}
  Exact Snell refraction for all sufficiently near-axis rays in the represented ray family. Angles are measured from the positive optical axis, so incidence and refraction angles relative to the surface normal are formed by subtraction. The ray family is allowed to have spherical aberration away from the axis; this predicate does not assert that all finite-height rays meet at PPrime
\end{definition}
\begin{proof}
  This is the declared model definition of Obeys Exact Snell Law Near Optical Axis.
\end{proof}
\begin{definition}[Has First Order Spherical Ray Geometry]
  \label{def:physics:phyx-mini-0041:phyxminiproblems-problemphyxmini0041-hasfirstordersphericalraygeometry}
  \lean{PhyXMiniProblems.ProblemPhyXMini0041.HasFirstOrderSphericalRayGeometry}
  \uses{def:physics:phyx-mini-0041:phyxminiproblems-problemphyxmini0041-lengthincentimeters, def:physics:phyx-mini-0041:phyxminiproblems-problemphyxmini0041-cylindricalglassrodsetup, def:physics:phyx-mini-0041:phyxminiproblems-problemphyxmini0041-objectdistancecentimeters, def:physics:phyx-mini-0041:phyxminiproblems-problemphyxmini0041-signedimagedistancecentimeters}
  The first-order geometry of the near-axis ray family. For surface height h, the incident direction has linear coefficient 1/s. The normal toward the center of curvature has coefficient -1/R. The refracted direction has coefficient -1/s', where s' is the signed paraxial axis intercept. These are local derivative statements, not global small-angle equalities and not the Gaussian surface equation itself
\end{definition}
\begin{proof}
  This is the declared model definition of Has First Order Spherical Ray Geometry.
\end{proof}
\begin{definition}[chief Ray Incident Angle Radians]
  \label{def:physics:phyx-mini-0041:phyxminiproblems-problemphyxmini0041-chiefrayincidentangleradians}
  \lean{PhyXMiniProblems.ProblemPhyXMini0041.chiefRayIncidentAngleRadians}
  \uses{def:physics:phyx-mini-0041:phyxminiproblems-problemphyxmini0041-cylindricalglassrodsetup, def:physics:phyx-mini-0041:phyxminiproblems-problemphyxmini0041-objectdistancecentimeters}
  Incident angle of the chief ray from an off-axis object to the vertex
\end{definition}
\begin{proof}
  This is the declared model definition of chief Ray Incident Angle Radians.
\end{proof}
\begin{definition}[chief Ray Refracted Angle Radians]
  \label{def:physics:phyx-mini-0041:phyxminiproblems-problemphyxmini0041-chiefrayrefractedangleradians}
  \lean{PhyXMiniProblems.ProblemPhyXMini0041.chiefRayRefractedAngleRadians}
  \uses{def:physics:phyx-mini-0041:phyxminiproblems-problemphyxmini0041-cylindricalglassrodsetup, def:physics:phyx-mini-0041:phyxminiproblems-problemphyxmini0041-signedimagedistancecentimeters}
  Refracted chief-ray angle from the vertex toward the signed paraxial image plane. For a virtual image, the ray is understood through backward extension
\end{definition}
\begin{proof}
  This is the declared model definition of chief Ray Refracted Angle Radians.
\end{proof}
\begin{definition}[Maps Optical Axis To Itself]
  \label{def:physics:phyx-mini-0041:phyxminiproblems-problemphyxmini0041-mapsopticalaxistoitself}
  \lean{PhyXMiniProblems.ProblemPhyXMini0041.MapsOpticalAxisToItself}
  \uses{def:physics:phyx-mini-0041:phyxminiproblems-problemphyxmini0041-cylindricalglassrodsetup}
  An on-axis object point maps to an on-axis image point
\end{definition}
\begin{proof}
  This is the declared model definition of Maps Optical Axis To Itself.
\end{proof}
\begin{definition}[Obeys Exact Chief Ray Snell Law Near Axis]
  \label{def:physics:phyx-mini-0041:phyxminiproblems-problemphyxmini0041-obeysexactchiefraysnelllawnearaxis}
  \lean{PhyXMiniProblems.ProblemPhyXMini0041.ObeysExactChiefRaySnellLawNearAxis}
  \uses{def:physics:phyx-mini-0041:phyxminiproblems-problemphyxmini0041-cylindricalglassrodsetup, def:physics:phyx-mini-0041:phyxminiproblems-problemphyxmini0041-chiefrayincidentangleradians, def:physics:phyx-mini-0041:phyxminiproblems-problemphyxmini0041-chiefrayrefractedangleradians}
  Exact Snell refraction at the surface vertex for chief rays from all sufficiently small object heights. The surface normal at the vertex is the optical axis, so no normal-angle correction appears
\end{definition}
\begin{proof}
  This is the declared model definition of Obeys Exact Chief Ray Snell Law Near Axis.
\end{proof}
\begin{definition}[Has Lateral Magnification At Optical Axis]
  \label{def:physics:phyx-mini-0041:phyxminiproblems-problemphyxmini0041-haslateralmagnificationatopticalaxis}
  \lean{PhyXMiniProblems.ProblemPhyXMini0041.HasLateralMagnificationAtOpticalAxis}
  \uses{def:physics:phyx-mini-0041:phyxminiproblems-problemphyxmini0041-cylindricalglassrodsetup}
  The physical definition of paraxial lateral magnification: the derivative of image height with respect to object height at the optical axis
\end{definition}
\begin{proof}
  This is the declared model definition of Has Lateral Magnification At Optical Axis.
\end{proof}
\begin{definition}[Answer Choice]
  \label{def:physics:phyx-mini-0041:phyxminiproblems-problemphyxmini0041-answerchoice}
  \lean{PhyXMiniProblems.ProblemPhyXMini0041.AnswerChoice}
  Labels of the four multiple-choice answers in the source problem
\end{definition}
\begin{proof}
  This is the declared model definition of Answer Choice.
\end{proof}
\begin{definition}[magnification Readout]
  \label{def:physics:phyx-mini-0041:phyxminiproblems-problemphyxmini0041-answerchoice-magnificationreadout}
  \lean{PhyXMiniProblems.ProblemPhyXMini0041.AnswerChoice.magnificationReadout}
  \uses{def:physics:phyx-mini-0041:phyxminiproblems-problemphyxmini0041-answerchoice}
  Dimensionless magnification printed for each answer choice
\end{definition}
\begin{proof}
  This is the declared model definition of magnification Readout.
\end{proof}
\begin{definition}[display Tolerance]
  \label{def:physics:phyx-mini-0041:phyxminiproblems-problemphyxmini0041-answerchoice-displaytolerance}
  \lean{PhyXMiniProblems.ProblemPhyXMini0041.AnswerChoice.displayTolerance}
  \uses{def:physics:phyx-mini-0041:phyxminiproblems-problemphyxmini0041-answerchoice}
  Half a unit in the last decimal place printed for an answer choice
\end{definition}
\begin{proof}
  This is the declared model definition of display Tolerance.
\end{proof}
\begin{definition}[Matches Answer Choice]
  \label{def:physics:phyx-mini-0041:phyxminiproblems-problemphyxmini0041-matchesanswerchoice}
  \lean{PhyXMiniProblems.ProblemPhyXMini0041.MatchesAnswerChoice}
  \uses{def:physics:phyx-mini-0041:phyxminiproblems-problemphyxmini0041-answerchoice}
  The exact result rounds to the readout displayed for the chosen answer
\end{definition}
\begin{proof}
  This is the declared model definition of Matches Answer Choice.
\end{proof}
\begin{lemma}[signed Image Distance Centimeters eq neg sixty four thirds]
  \label{lem:physics:phyx-mini-0041:phyxminiproblems-problemphyxmini0041-signedimagedistancecentimeters-eq-neg-sixty-four-thirds}
  \lean{PhyXMiniProblems.ProblemPhyXMini0041.signedImageDistanceCentimeters_eq_neg_sixty_four_thirds}
  \uses{def:physics:phyx-mini-0041:phyxminiproblems-problemphyxmini0041-cylindricalglassrodsetup, def:physics:phyx-mini-0041:phyxminiproblems-problemphyxmini0041-signedimagedistancecentimeters, def:physics:phyx-mini-0041:phyxminiproblems-problemphyxmini0041-hasstatedrefractiveindices, def:physics:phyx-mini-0041:phyxminiproblems-problemphyxmini0041-matchesfigurereadout, def:physics:phyx-mini-0041:phyxminiproblems-problemphyxmini0041-obeysexactsnelllawnearopticalaxis, def:physics:phyx-mini-0041:phyxminiproblems-problemphyxmini0041-hasfirstordersphericalraygeometry}
  Differentiating exact Snell refraction at the optical axis, using the three first-order geometric relations, locates the paraxial virtual image at s' = -64/3 cm. This is an intermediate consequence, not a figure readout
\end{lemma}
\begin{proof}
  The conclusion follows from the governing relations and figure readouts named in the statement of signed Image Distance Centimeters eq neg sixty four thirds.
\end{proof}
% --- Archon declaration topology end ---
% --- Archon physics formalization source end ---
